%% notes-tex/multimodal/043-training-plan/sections/6-launch.tex
%% SOURCE: .claude/skills/prepare-experiment/SKILL.md (sections 1 to 5) at commit bed26f949;
%% experiments/043-expression-proteome-morphology/scripts/wandb_view_043.py and
%% results/wandb_view_ids.json; memories delta-never-preempt-partition,
%% always-show-train-and-validation-loss-together, wandb-views-all-runs-smoothing-loss-first.

\section{Launch discipline\secstatus{todo}}
\label{sec:launch}

\notesec{The gate}
Every machine's first submission is a smoke, never a chain: G1 on GilaHyper, P0 on Polaris,
A0 and A1 on Aurora, D7 on Delta. After a smoke has logged two evaluations,
\file{scripts/check_wandb_run_keys.py} is run on its W\&B run with \texttt{-{}-require} for
\texttt{val/loss} and, for every dataset of the arm, \texttt{val/<dataset>/loss} and its
validation score; exit 1 means the trainer is fixed and the smoke rerun before anything depends
on it. The gate is rerun after every trainer change, and its output (key, row count, verdict)
goes verbatim into the round's dated section of the experiment note.

\notesec{The logging contract}
At every evaluation, as curves: \texttt{epoch}, \texttt{train/loss} and \texttt{val/loss} under
the same weighting, the per-dataset pair \texttt{train/<dataset>/loss} and
\texttt{val/<dataset>/loss}, \texttt{train/lr}, \texttt{train/grad\_norm}, every validation
statistic of Table~\ref{tab:ladder}, and \texttt{operator/relative\_norm}, the operator's
attention branch over the context norm. A stage-2 run also logs \texttt{stage}, the stage-1 run id it started
from and the checkpoint's sha256.

\notesec{Time}
\texttt{-{}-time} and PBS walltime are about 1.5 times the estimate from a measured step time on
that machine (Sec.~\ref{sec:cost}), never a round number chosen in advance. A step-matched run
sets \texttt{budget.max\_steps}, and the trainer raises rather than train a shorter unmatched
control.

\notesec{Chains}
Slurm chains use \texttt{-{}-dependency=afterany}, PBS chains \texttt{-W depend=afterany:<id>},
so a crash never stalls a chain; a fine-tune whose checkpoint is missing raises at startup.
GilaHyper uses at most the two cards granted. Delta never uses the preempt partition and never
launches without the owner's go. Polaris uses no \texttt{preemptable} queue while the trainer
has no resume.

\notesec{W\&B}
One project, \texttt{torchcell\_043-expression-proteome-morphology}. Groups are named
\file{s1_<arm>_<loss>_<weighting>} and \file{s2_<target>_<from-arm>_<lr>_<policy>},
runs \file{<group>_s<seed>}, with the machine as a tag. \file{wandb_view_043.py} writes the
round's view before the first result is read, loss first (train and validation on one panel,
grouped by arm, smoothing on, with a step-axis twin where stage-2 runs and their controls differ
in epochs per step), then the headline validation statistics, the per-dataset panels, the train
side and bookkeeping. Aborted runs are tagged, not deleted. Links go one per line.

\notesec{What is committed where}
Configs (\file{conf/arm/all.yaml}, \file{conf/stage2/*.yaml}, \file{conf/decision_rule.yaml}),
launchers (\file{scripts/gh_*.slurm}, \file{scripts/delta_*.slurm}, \file{scripts/polaris_*.pbs},
\file{scripts/aurora_*.pbs}) and the task lists of the packed jobs live in
\file{experiments/043-expression-proteome-morphology/} and are committed before submission.
Window means, ladder reads and probe reads are JSON files under \file{results/}, each written by
a committed script from the unsampled W\&B history or the checkpoints. Each round gets a dated
section in the experiment note with the gate output, the queue table (job id, run, dependency,
time requested, expected end) and the arms dropped for budget, by name. Checkpoints, prepared
stores and job logs stay under \texttt{\$DATA\_ROOT} and the \texttt{/eagle} and \texttt{/flare}
project trees, never in git. The numbers of this document are hand-sourced with
\texttt{\%\% SOURCE} comments; a table script in the experiment folder replaces them when the
first stage-1 results arrive.
